% !TeX encoding = UTF-8
% !TeX root = thesis.tex

\section{Fazit und Ausblick}

% ---------------------------------------------------------------------------
\subsection{Zusammenfassung}
% ---------------------------------------------------------------------------

Diese Arbeit hat eine zweistufige hybride Pipeline zur Reduzierung von
False Positives in der Cloud-Anomalieerkennung entwickelt und auf dem
synthetischen Benchmark CloudAnoBench evaluiert. Der methodische Fokus lag
auf einer leakage-sicheren Evaluation unter einer festen
Recall-Nebenbedingung ($R \geq 0{,}95$), nicht auf der Maximierung
einzelner Metriken. Beide Komponenten~--- metrischer Detektor und
Log-Verifikation~--- wurden isoliert und in zwei Fusionskonfigurationen
evaluiert; alle Ergebnisse sind durch einen festen Zufallszustand und
13~Unit-Tests reproduzierbar.

% ---------------------------------------------------------------------------
\subsection{Beantwortung der Forschungsfrage}
% ---------------------------------------------------------------------------

Die Forschungsfrage wird bejaht: Die Log-Komponente reduziert die FPR von
$79{,}8\,\%$ auf $26{,}3\,\%$ bei einem Recall von $0{,}958$~--- die
Recall-Nebenbedingung ist damit erfüllt. Die vollständige Ergebnisanalyse
und Argumentation findet sich in Abschnitt~\ref{sec:ablation} und den
Abschnitten zur Diskussion.

Der qualifizierende Befund ist, dass Log-Kontext auf Case-Ebene als
eigenständiger Prädiktor dominiert und die Fusionsstufe keinen zusätzlichen
Beitrag liefert: Das AND-Gate verliert aufgrund des niedrigen
Metrik-Thresholds seine Filterfunktion auf Metrikseite; die Soft-Fusion
erzielt eine niedrigere AP als die Log-Komponente allein. Die Nützlichkeit
einer AND-Fusion hängt damit entscheidend von der Kalibrierung des
Metrik-Triggers auf Case-Ebene ab, die in dieser Arbeit nicht umgesetzt wurde.

Diese Befunde gelten für den CloudAnoBench-Datensatz unter den gewählten
Konfigurationen und sind ohne weitere Evidenz nicht zu verallgemeinern.

% ---------------------------------------------------------------------------
\subsection{Ausblick}
% ---------------------------------------------------------------------------

Aus den Befunden dieser Arbeit ergeben sich mehrere konkrete Ansatzpunkte
für zukünftige Arbeiten:

\begin{description}
  \item[Case-Level-Kalibrierung des Metrik-Triggers.]
    Die AND-Gate-Degeneration ist eine direkte Folge der Threshold-Wahl auf
    Zeilenebene. Eine Kalibrierung direkt auf aggregierten Case-Scores~---
    oder der Einsatz alternativer Aggregationsregeln wie Median statt
    Maximum~--- würde die AND-Gate-Logik erst operativ wirksam machen und
    den eigenständigen Beitrag der Metrik-Komponente zur Fusion prüfbar machen.

  \item[Alternative Fusionsregeln.]
    Neben AND-Gate und Produkt-Regel existieren weitere Fusionsstrategien~---
    etwa Summen-Regel, gewichtete Mehrheitsentscheidung oder lernbasierte
    Kombinatoren. Ob solche Regeln auf CloudAnoBench einen messbaren Vorteil
    gegenüber der Log-Komponente allein erzielen, bleibt eine offene Frage.

  \item[Weiterentwicklung der Log-Komponente.]
    Die log-basierte Verifikation wurde in dieser Arbeit als interpretierbare
    Baseline mit TF-IDF und logistischer Regression umgesetzt. Ob
    struktur-bewusste Log-Parser, vortrainierte Sprachmodelle oder
    Embedding-Methoden auf diesem Datensatz zu einer weiteren FPR-Reduktion
    führen, wäre eine naheliegende Erweiterung.

  \item[Statistische Absicherung.]
    Die berichteten Metriken basieren auf einem einzigen Trainings"=Test"=Split.
    Die Frage, wie stabil die Ergebnisse über verschiedene Splits sind, ließe
    sich durch Bootstrap-Konfidenzintervalle oder wiederholte Splits mit
    unterschiedlichen Zufallszuständen beantworten.

  \item[Evaluation auf realem Produktions-Traffic.]
    CloudAnoBench verwendet kontrollierte Anomalie-Injektionen. Ob die
    gelernten Muster auf reale Infrastrukturdaten~--- mit natürlichen
    Klassenverteilungen, variablen Log-Formaten und unbekannten
    Anomalietypen~--- übertragbar sind, ist durch diese Arbeit nicht
    untersucht und stellt eine wichtige Anschlussfrage dar.

  \item[Modularisierung der Pipeline.]
    Die funktionale Pipeline liegt derzeit überwiegend in Jupyter Notebooks.
    Eine Modularisierung in eine paketierbare, konfigurierbare
    Software-Pipeline würde die Wiederverwendbarkeit und Anpassbarkeit an
    andere Datensätze wesentlich verbessern.
\end{description}
